% !TEX TS-program = pdflatex
% !TEX encoding = UTF-8 Unicode

% This is a simple template for a LaTeX document using the "article" class.
% See "book", "report", "letter" for other types of document.

\documentclass[11pt]{article} % use larger type; default would be 10pt

\usepackage[utf8]{inputenc} % set input encoding (not needed with XeLaTeX)

%%% Examples of Article customizations
% These packages are optional, depending whether you want the features they provide.
% See the LaTeX Companion or other references for full information.

%%% PAGE DIMENSIONS
\usepackage{geometry} % to change the page dimensions
\geometry{a4paper} % or letterpaper (US) or a5paper or....
% \geometry{margin=2in} % for example, change the margins to 2 inches all round
% \geometry{landscape} % set up the page for landscape
%   read geometry.pdf for detailed page layout information

\usepackage{graphicx} % support the \includegraphics command and options

% \usepackage[parfill]{parskip} % Activate to begin paragraphs with an empty line rather than an indent

%%% PACKAGES
\usepackage{booktabs} % for much better looking tables
\usepackage{array} % for better arrays (eg matrices) in maths
%\usepackage{paralist} % very flexible & customisable lists (eg. enumerate/itemize, etc.)
\usepackage{verbatim} % adds environment for commenting out blocks of text & for better verbatim
\usepackage{subfig} % make it possible to include more than one captioned figure/table in a single float
% These packages are all incorporated in the memoir class to one degree or another...

%%% HEADERS & FOOTERS
\usepackage{fancyhdr} % This should be set AFTER setting up the page geometry
\pagestyle{fancy} % options: empty , plain , fancy
\renewcommand{\headrulewidth}{0pt} % customise the layout...
\lhead{}\chead{}\rhead{}
\lfoot{}\cfoot{\thepage}\rfoot{}

%%% SECTION TITLE APPEARANCE
\usepackage{sectsty}
\allsectionsfont{\sffamily\mdseries\upshape} % (See the fntguide.pdf for font help)
% (This matches ConTeXt defaults)

%%% ToC (table of contents) APPEARANCE
\usepackage[nottoc,notlof,notlot]{tocbibind} % Put the bibliography in the ToC
\usepackage[titles,subfigure]{tocloft} % Alter the style of the Table of Contents
\renewcommand{\cftsecfont}{\rmfamily\mdseries\upshape}
\renewcommand{\cftsecpagefont}{\rmfamily\mdseries\upshape} % No bold!

%%% END Article customizations

\usepackage[spanish]{babel}
\usepackage{listings} 
%%% The "real" document content comes below...

\title{Python}

%\date{} % Activate to display a given date or no date (if empty),
         % otherwise the current date is printed 

\begin{document}
\maketitle

%\tableofcontents % No hace falta un TOC en un artículo corto

\setlength{\unitlength}{1 cm} %Especificar unidad de trabajo
\thispagestyle{empty}
\begin{picture}(18,4)
\put(0,0){\includegraphics[width=3cm,height=4cm]{imagenes/logoEspol.jpg}}
\put(11.5,0){\includegraphics[width=4cm,height=4cm]{imagenes/logoFiec.jpg}}
\end{picture}
\\
\\
\begin{center}
\textbf{{\Huge Lenguajes de Programación}\\[0.5cm]
{\LARGE Proyecto de Python}}\\[1.25cm]
{\LARGE \textbf{Grupo DJM}}\\[1cm]
{\LARGE \textbf{ Integrantes:}}\\[1cm]
{\Large Denisse Pintado}\\[1cm]
{\Large Jonathan Mendieta}\\[1cm]
{\Large Janina Costa}\\[1cm]

\newpage
\end{center}

\newpage
\title{Python}

\section{Introducción}
El siguiente proyecto es sobre el lenguaje de programación Python.\\ 
\\En el siguiente documentos redactaremos detalladamente el funcionamiento de nuestro código, el procedamiento y los pasos que utilizamos para desarrollar el codigo \\
\\Para conocer más sobre este lenguaje de programación veremos las características más importantes,  como surgió y los requisitos para su correcta instalación. \\
\\El proyecto se realiza por el interés de aprender sobre un nuevo lenguaje el cual nos aporte más conocimientos sobre el amplio mundo de la programación.\\
\\
\\ \\

OBJETIVOS:
\\ \\
\\El objetivo es poder leer de un archivo XML.\\\\\\

\section{Python}
Python - Ventajas y Desventajas
Python es un lenguaje interpretado, orientado a objetos de propósito general. Python permite mantener de forma sencilla interacción con el sistema operativo, y resulta muy adecuado para manipular archivos de texto. 
\newline
\\
\newline
\\


\section{ Ventajas y Desventajas}
VENTAJAS\\
Desarrollo más rápido : Puedes escribir un programa, salvarlo y ejecutarlo. En un lenguaje compilado tienes que pasar por los pasos de compilar y ligar el software, lo cual puede ser un proceso lento.
Multiplataforma : El mismo código funciona en cualquier arquitectura, la única condición es que disponga del intérprete del lenguaje. No es necesario compilar el código una vez para cada arquitectura.\\\\
\newline
DESVENTAJAS\\\\
Lentitud : Los programas interpretados son más lentos que los compilados. Sin embargo los programas interpretados suelen ser cortos, en los que la diferencia es inapreciable.
\newline
\\
\newline
\\

\section{ Problemas que se presentaron}





\end{document}